\section{Creating a local repository\label{sec:local_repo}}

\CNAME\ is stored on \href{http://www.github.com/Casal2/}{GitHub} --- and uses it to manage and store copies of the source code and ancillary files. This section is intended to be an introduction for new users to GitHub and git, and covers the following:

\begin{enumerate}
	\item Registering a GitHub username
	\item Using git
	\item Cloning a repository to your local machine
	\item Fork the master repository
\end{enumerate}

\subsection{Git username and profile}

The first step is to create a username and profile on GitHub (if you do not already have one). Go to \url{https://www.github.com}, register a username, and set up a profile if you do not already have one. See the GitHub help for more information. 

Once you have set up a GitHub account, you need to download git software (see below) so you can access and modify the code on your local machine.

\subsection{Using git}

You will need to install and use git software in order to download ('clone') the repository to your local machine and upload ('push') changes to your local repository. There are several git clients available and we suggest \href{https://desktop.github.com/}{GitHub Desktop} on Microsoft Windows as a good starting point. 

\CNAME\ also requires a command line version of git in order to compile (see Section~\ref{sec:build_environment} for a complete description of the requirements of the \CNAME\ build environment).  

Download a git client from \href{https://git-scm.com/downloads}{git software} and you ensure that the binary is included into your system path. This can be checked by opening a terminal (e.g., Powershell or command prompt on Microsoft Windows machines and typing in \texttt{git -v}. 

\subsection{Cloning a repository to your local machine}\label{sec:CloningGitHub}

The \CNAME\ code, documentation, and other ancillary files is kept in the Master branch of the repository. Only the \CNAME\ Development Team have permission to add, delete, or change code directly. You can clone this repository to your local machine, but note that you won't be able to 'push' any of your changes back to the \CNAME\ Master repository. See the below in Section \ref{sec:maintain_repo} for creating your own copy ('fork') of \CNAME\ on GitHub.

To clone the repository, navigate to \url{https://github.com/Casal2/CASAL2} and use the green code button, which is shown below 'clone'. There are multiple protocols (https, ssh etc.) for cloning highlighted in the red box. If you have GitHub Desktop installed, then choose this option and this will open the link in GitHub Desktop on your local machine (see the circled text in Figure~\ref{fig:clone}).

\begin{figure}[H]
	\centering
	\includegraphics[scale=0.6]{Figures/clone_repo.png}
	\caption{Cloning a repository}\label{fig:clone}
\end{figure}

\subsection{Create a GitHub fork to start making changes\label{sec:maintain_repo}}

This section covers how to fork a copy of the Master codebase of \CNAME\ into your GitHub account.  

If you wish to contribute to \CNAME, you will need to fork a copy of \CNAME\ to your GitHub account. This provides a way for you to have a copy of the code and upload changes for your local version. Once you have finalised a change, it can then be submitted to the \CNAME\ Development Team for evaluation and potential inclusion into the main code base. 

To do this, simply navigate to the \CNAME\ GitHub repository, and select fork, and then create a new fork. 

As described above in Section~\ref{sec:CloningGitHub}, this fork can be cloned to your local machine and used to modify, add, or delete code. Note that any changes made here only effect your fork and will not modify the Master codebase at \url{https://github.com/Casal2/CASAL2}.

Once you have made a push, you will then be able to see the changes made. If you have any difficulty, check the GitHub documentation for help.


\subsection{Adding an SSH authentication key to your account on GitHub.com}

One of the reasons contributors may have difficulty pushing and pulling changes to \CNAME\ is because they don't have a SSH authentication key linked to their GitHub account. The error message may look like that shown in Figure~\ref{fig:permissiondenied}. Please see the information from \href{https://docs.github.com/en/authentication/connecting-to-github-with-ssh/adding-a-new-ssh-key-to-your-github-account}{this link} for detailed instructions on how to create a public SSH key on your computer and then link it to your GitHub account.

\begin{figure}[H]
	\centering
	\includegraphics[scale=0.6]{Figures/permissiondenied.png}
	\caption{Example of the error message when SSH authentication has not been enabled}\label{fig:permissiondenied}
\end{figure}
